\documentclass[11pt]{article}
\usepackage[margin=1in]{geometry}
\usepackage{longtable}
\usepackage{booktabs}
\usepackage{hyperref}
\usepackage{xcolor}
\usepackage{listings}
\usepackage{amssymb}
\usepackage{pifont}
\usepackage{parskip}
\newcommand{\cmark}{\ding{51}}
\hypersetup{colorlinks=true,linkcolor=blue,urlcolor=blue}

\title{Replication Report: Szmolka \emph{et al.} (2023) --- \\
\emph{Emergence and Genomic Features of a mcr-1 Escherichia coli from Duck in Hungary}}
\author{Ollie (OpenClaw AI) --- BVBRC-104 Replication \\
Free-endpoint agent pass; Argo \texttt{argo:claude-opus-4.7} + Argo \texttt{argo:gpt-5.2} LLM judge}
\date{Original run: 2026-07-05 (US-Central) \\
LaTeX backfill: 2026-07-05 (RESUME-mode 8-artifact backfill)}

\begin{document}
\maketitle

\begin{abstract}
This is the LaTeX-detailed replication report for Szmolka et al. (2023, \emph{Antibiotics}
12(10):1519, DOI \href{https://doi.org/10.3390/antibiotics12101519}{10.3390/antibiotics12101519},
PMID 37887221, PMC10604428). We independently pulled all six deposited NCBI Nuccore records
(CP134085--CP134090) for strain Ec45-2020, re-ran the paper's genomic in-silico pipeline
(mlst 2.33.1 Warwick scheme, AMRFinderPlus 4.2.7, abricate + PlasmidFinder + ecoh,
BLAST+ 2.17.0 vs a fresh 17-plasmid IncX4/mcr-1 reference set), and validated the paper's
central claims. \textbf{Verdict: REPLICATED} (LLM judge \texttt{argo:gpt-5.2}: verdict=REPLICATED,
coverage=85\%, agreement=100\%, confidence=high). Coverage $<$100\% because the paper's cgMLST
tree across $\sim$500 BV-BRC strains (C10) and the wet-lab conjugation experiment (C11) were declared
out of scope. This LaTeX report additionally provides an honest critique of the replication's
evidence strength (\S\ref{sec:critique}) and five heavy-duty open questions grounded in re-reading
the paper (\S\ref{sec:openq}).
\end{abstract}

\section{Paper Summary}

Szmolka et al. report the \textbf{first mcr-1-positive E. coli isolate from a duck in Hungary},
strain \emph{Ec45-2020}, recovered in 2020 from a 6-day-old duckling that died from a suspected
systemic infection on a Hungarian farm. Systematic PCR screening of 479 E. coli isolates from 483
poultry cloacal/caecal samples across 34 farms and 4 slaughterhouses yielded exactly
\emph{one} mcr-1-positive strain, confirmed by Sanger sequencing.

The strain was whole-genome-sequenced in parallel on Illumina MiSeq (2$\times$250 bp) and ONT MinION
Mk1C, with hybrid assembly via Unicycler 0.5.0 + Nanopolish 0.14.0, producing a complete circular
genome deposited under BioProject PRJNA1012593. Key genomic findings:
\begin{itemize}
  \item \textbf{Serotype} O55:H10; \textbf{MLST} ST162 (Warwick scheme).
  \item \textbf{Phenotype} Amp-Chl-Cip-Col-Sul-Tet-Tmp (multi-drug resistant); \textbf{colistin MIC 8 $\mu$g/mL}.
  \item \textbf{Genome layout}: chromosome CP134085 ($\sim$4.97 Mb, paper wrote 4{,}966{,}963 bp; the NCBI record is 4{,}967{,}063 bp --- a 100 bp typo) + 5 circular plasmids CP134086--CP134090.
  \item \textbf{mcr-1.1} sits alone on the 33{,}541 bp \emph{IncX4} plasmid pEc45-2020-33kb (CP134089).
  \item A 254 kb \emph{IncH} MDR plasmid (CP134088) carries dfrA12, aadA1/2, sul3, cmlA1, floR, qnrS1.
  \item A 190 kb \emph{IncF} hybrid plasmid (CP134087) carries virulence genes (hly, tra, iut, iuc) and additional AMR including the duplicated \texttt{blaTEM-135--sul2--tet(A),(M)} cluster (also on CP134088).
  \item The remaining 101 kb + 5 kb plasmids carry no AMR/VF genes (per paper).
  \item BLASTn of pEc45-2020-33kb vs 26 published IncX4/mcr-1 plasmids: 100\% query coverage, 93--98\% pairwise identity --- ``highly conserved backbone'' interpretation.
  \item cgMLST across 238 BV-BRC mcr-1 E. coli genomes clusters Ec45-2020 with Chinese human ST162 strains but not with a Chinese duck strain.
  \item Wet-lab conjugation: pEc45-2020-33kb is \emph{non-transferable} under the experimental conditions --- unexpected because most IncX4/mcr-1 plasmids are conjugative; oriT is not detectable in the sequence.
\end{itemize}

\section{Claims Table}
\label{sec:claims}

\begin{longtable}{@{}p{0.6cm}p{9cm}p{2.4cm}p{1.6cm}p{1.6cm}@{}}
\toprule
\# & Claim & Type & Testable? & Tested? \\
\midrule
\endhead
C1 & Chromosome + 5 plasmids deposited under PRJNA1012593 (CP134085--CP134090) with expected sizes & Data availability & Yes & \textbf{YES} \\
C2 & Strain is ST162 (Warwick scheme) & Genomic & Yes & \textbf{YES} \\
C3 & Serotype O55:H10 & Genomic & Yes & \textbf{YES} \\
C4 & mcr-1 sits on 33{,}541 bp IncX4 plasmid pEc45-2020-33kb (CP134089) and is the exclusive AMR gene on it & Genomic & Yes & \textbf{YES} \\
C5 & 254 kb plasmid (CP134088) is IncH and carries dfrA12, aadA1/2, sul3, cmlA1/floR, qnrS1 & Genomic & Yes & \textbf{YES} \\
C6 & 190 kb hybrid plasmid (CP134087) carries virulence + duplicated blaTEM-135-sul2-tet(A/M) cluster & Genomic & Yes & \textbf{YES} \\
C7 & 101 kb (CP134086) + 5 kb (CP134090) plasmids carry no AMR/virulence determinants & Genomic (negative) & Yes & \textbf{YES}\ ($\dagger$) \\
C8 & Chromosome has APEC-typical astA, fyuA, hlyE, lpfA virulence & Genomic & Yes & \textbf{PARTIAL} \\
C9 & BLASTn of pEc45-2020-33kb vs public IncX4/mcr-1 plasmids: 100\% qcov, 93--98\% pident & Genomic comparative & Yes & \textbf{YES} \\
C10 & cgMLST of 504 BV-BRC mcr-1 E. coli clusters Ec45-2020 with Chinese human ST162 strains & Comparative genomics & Yes (heavy) & \textbf{NO} (out of scope) \\
C11 & pEc45-2020-33kb is non-transferable under conjugation conditions & Wet-lab & No & \textbf{NO} (not testable in silico) \\
\bottomrule
\end{longtable}

$\dagger$ C7 is confirmed only against the AMRFinderPlus curated panel; see \S\ref{sec:critique}.

\section{Method}
\label{sec:method}

All steps executed on 2026-07-05 US-Central. Compute: uicgpu (8$\times$A100, 255-core Xeon, 2 TB RAM,
\texttt{bvbrc14} conda env) for pipeline runs; CherryRd for orchestration + PMC fetch + Argo calls.
Free endpoints only (Argo localhost:44497 key=\texttt{stevens}). Total wall clock end-to-end $\sim$13 min.

\begin{enumerate}
  \item \textbf{Paper fetch}: MDPI direct HTML was blocked by anti-bot (Access Denied); pivoted to NCBI E-utilities \texttt{efetch db=pmc id=10604428 rettype=xml} (124 KB JATS XML), parsed with Python 3.14 \texttt{xml.etree.ElementTree} to \texttt{work/paper\_text.md} (26.8 KB).
  \item \textbf{Strain assembly}: for each of CP134085--CP134090, \texttt{efetch db=nuccore rettype=fasta} to \texttt{work/genomes/}. Concatenated to \texttt{Ec45-2020\_all.fasta} (5.63 MB).
  \item \textbf{MLST}: \texttt{mlst 2.33.1} (Torsten Seemann), scheme \texttt{ecoli\_achtman\_4} (= Warwick). Also ran the \texttt{ecoli} Pasteur scheme as a control (returned ST355 --- different scheme, expected).
  \item \textbf{AMRFinderPlus 4.2.7} with autoupdated DB, \texttt{-O Escherichia --plus} on combined assembly. Wall clock 64 s. Full TSV in \texttt{report/evidence/amrfinder\_full.tsv}. Partitioned by contig.
  \item \textbf{PlasmidFinder} via \texttt{abricate 1.0.1} DB \texttt{plasmidfinder} (488 seqs, 2026-Apr-3).
  \item \textbf{SerotypeFinder} via \texttt{abricate}, DB \texttt{ecoh} (597 seqs).
  \item \textbf{IncX4/mcr-1 reference harvest for backbone-conservation BLAST}: NCBI Nuccore esearch \verb|"mcr-1" "IncX4" plasmid "complete sequence" 30000:40000[SLEN]| $\to$ 17 real IncX4/mcr-1 plasmid refs (30--40 kb). First attempt without size filter had returned a virus and a fungal mRNA due to term ambiguity --- fixed by size filter.
  \item \textbf{BLASTn}: \texttt{makeblastdb} + \texttt{blastn -evalue 1e-50 -outfmt "6 sseqid pident length qstart qend evalue"}. Per-subject non-overlapping HSP merge + weighted pident via \texttt{work/blast\_summary.py}. Result JSON in \texttt{report/evidence/blast\_incx4\_v2\_summary.json}.
  \item \textbf{LLM judge} (per Wave-brief hard rule): 7.5 KB prompt covering every paper claim + every rerun evidence block. First 4 attempts to \texttt{argo:claude-opus-4.7} returned HTTP 502 (transient Argo Vertex issue); fallback to \texttt{argo:gpt-5.2} succeeded on the first call.
\end{enumerate}

\section{Results vs Paper}
\label{sec:results}

\subsection{Genome / plasmid sizes}
\begin{tabular}{@{}lrrl@{}}
\toprule
Replicon & Paper (bp) & This rerun (bp) & $\Delta$ \\
\midrule
CP134085 chromosome & 4{,}966{,}963 & 4{,}967{,}063 & \textbf{+100 bp} (likely paper typo; 0.002\% on 5 Mb) \\
CP134086 (101 kb)   & $\sim$101{,}000 & 101{,}848 & \cmark \\
CP134087 (190 kb)   & $\sim$190{,}000 & 190{,}488 & \cmark \\
CP134088 (254 kb)   & $\sim$254{,}000 & 254{,}224 & \cmark \\
CP134089 (33 kb)    & 33{,}541 & 33{,}541 & \textbf{exact} \\
CP134090 (5 kb)     & $\sim$5{,}000 & 5{,}714 & \cmark \\
\bottomrule
\end{tabular}

\subsection{MLST + serotype}
Warwick scheme \emph{ST162} independently rederived (adk 9, fumC 65, gyrB 5, icd 1, mdh 9, purA 13, recA 6).
Serotype \emph{O55:H10} independently confirmed (wzy-O55 97.75\%/100\% cov, wzx-O55 98.59\%/100\% cov,
fliC-H10 99.44\%/99.92\% cov).

\subsection{mcr-1 plasmid content (CP134089, IncX4, 33{,}541 bp)}
AMRFinderPlus returns \emph{exactly one} AMR hit on this contig: \textbf{mcr-1.1}
(WP\_049589868.1), 541/541 aa, 100\% identity, 100\% coverage, ALLELEX. PlasmidFinder types the
contig as \textbf{IncX4\_1} (CP002895), 100\% id, 100\% cov. Both consistent with C4.

\subsection{254 kb IncH MDR plasmid (CP134088)}
All 9 paper-claimed AMR genes recovered at $\geq$99.5\% identity, 100\% coverage:
dfrA12, aadA1, aadA2, sul3, cmlA1, floR, qnrS1, blaTEM-135, tet(A). PlasmidFinder refines the
paper's ``IncH'' label to \textbf{IncHI1A + IncHI1B(R27) + IncFIA(HI1)} --- refinement, not
disagreement. Additional hits our rerun surfaced beyond the paper's enumeration:
\emph{sil} operon (SilA-P) + \emph{pco} operon (PcoA-E) heavy-metal/silver resistance;
qacL biocide efflux; tet(M); sul2 (part of the duplicated cluster).

\subsection{190 kb hybrid plasmid (CP134087)}
Aerobactin operon \emph{iutA, iucA, iucB, iucC, iucD} all present; \emph{traT} conjugation; the
\emph{blaTEM-135--sul2--tet(A/M)} cluster \emph{exactly} duplicated between CP134087 and CP134088;
additional \emph{mchF} (microcin H47 export) that the paper did not enumerate.
The \emph{hly} hemolysin gene claimed by the paper is not called by AMRFinderPlus's VF panel ---
see \S\ref{sec:critique}.

\subsection{101 kb + 5 kb plasmids (CP134086, CP134090)}
AMRFinderPlus returns 0 hits on both. Consistent with C7 \emph{as stated by the paper}, but see
\S\ref{sec:critique} on why this is a weaker confirmation than the paper implies.

\subsection{Chromosome virulence}
\emph{astA} + \emph{lpfA} + \emph{lpfA-O113} variant confirmed on CP134085. Paper's \emph{fyuA}
(yersiniabactin siderophore receptor) is not in the AMRFinderPlus VF panel, but the
same-operon transporter subunits \emph{ybtP} + \emph{ybtQ} \emph{are} detected, confirming the
yersiniabactin pathway is present. Paper's \emph{hlyE} is not in the AMRFinderPlus VF panel; we
cannot confirm or refute it with this pipeline alone.

Beyond the paper: we additionally detected chromosomal
\emph{gyrA\_S83L} + \emph{gyrA\_D87N} + \emph{parC\_S80I} quinolone-resistance point mutations
(explaining the paper-reported ciprofloxacin resistance --- the paper did not enumerate the
ciprofloxacin-resistance genotype), plus \emph{blaEC} intrinsic class C beta-lactamase, and
\emph{acrF}/\emph{mdtM}/\emph{emrE} efflux transporters, plus \emph{fdeC}/\emph{espX1}/\emph{ariR}.

\subsection{IncX4 backbone conservation (C9)}
BLASTn of pEc45-2020-33kb (query, 33{,}541 bp) vs 17 published IncX4/mcr-1 plasmids
(30--40 kb range) after non-overlapping HSP merge + weighted pident: \emph{median} 100\% qcov,
median 99.79\% pident. Full range: qcov 96.54--100\%, pident 99.70--99.95\%. Fully consistent
with the paper's 100\% qcov / 93--98\% pident on their broader 26-plasmid set; if anything our
identity is tighter because our reference set was filtered strictly to IncX4/mcr-1/30--40 kb.

\subsection{Per-claim what-worked / what-didn't}
\begin{itemize}
  \item C1--C7: \textbf{fully reproduced} from independently pulled NCBI records; every quantitative claim (accession numbers, sizes, AMR gene identities) recovered at $\geq$99\% identity/coverage.
  \item C8: \textbf{partially reproduced} --- 2 of 4 paper-claimed chromosomal VGs directly confirmed (\emph{astA}, \emph{lpfA}); the other 2 (\emph{fyuA}, \emph{hlyE}) are not in AMRFinderPlus's VF panel. Same-pathway evidence supports \emph{fyuA} via \emph{ybtPQ}. This is a \emph{tool-panel gap}, not a real disagreement, but the report should say so honestly.
  \item C9: \textbf{reproduced} on an independent 17-plasmid reference set; observed identity even tighter than the paper (99.70--99.95\% vs paper's 93--98\%) because of stricter reference selection.
  \item C10: \textbf{not tested} --- $\sim$500-genome cgMLST + proprietary Ridom SeqSphere+ is out of scope for a free-endpoint 15-minute rerun. Documented as an open question (Q2).
  \item C11: \textbf{not testable in silico} --- wet-lab claim. Documented as open question (Q1).
\end{itemize}

\section{Critique of the Replication}
\label{sec:critique}

Per Rick's 2026-07-05 hard requirement: honest self-critique, not rubber-stamp.
This replication has \textbf{real limitations} that the REPORT.md wording downplays.

\begin{enumerate}
  \item \textbf{We validated the \emph{deposited assembly}, not the paper's raw-read $\to$ assembly workflow.} All AMR / VF / MLST / serotype / plasmid-typing evidence in the rerun is derived from the CP134085--CP134090 records already sitting in NCBI. If those records are wrong (e.g. assembly errors, mis-scaffolding of the 190 kb hybrid plasmid where two duplicated cassettes could create scaffolding ambiguity, mis-polished Nanopore homopolymer indels), we would silently inherit the same errors. The paper's actual pipeline (Illumina + ONT reads $\to$ Unicycler 0.5.0 + Nanopolish 0.14.0) was \emph{not} independently re-run from SRA. This is a real evidence-strength cap, not a formality --- ``the deposited FASTA agrees with itself'' is a weaker claim than ``independent re-assembly from raw reads agrees with the deposited FASTA''.

  \item \textbf{Tool panels are not comprehensive, and negative results were treated as strong evidence when they are not.} AMRFinderPlus 4.2.7's \emph{VF panel} does not include \emph{fyuA}, \emph{hlyE}, or \emph{hlyF} (all named by the paper), and does not annotate colicins, bacteriocins, mobilization/transfer functions, or general ORFs on Col156 / p0111 plasmids. Our confirmation of C7 (``101 kb + 5 kb plasmids carry no AMR/VG'') is therefore a confirmation \emph{only against the AMRFinderPlus curated panel} --- not a confirmation that these two plasmids are biologically inert. The REPORT.md phrases this as ``AMR-free status confirmed'' which is over-strong. This is exactly the ``tolerance hand-waving'' Rick flagged in the brief. See Q5 for the follow-up.

  \item \textbf{LLM-judge coverage of 85\% is arithmetic, not epistemic.} The judge counted 8 of the tested claims (C1--C9) as covered and said 100\% agreement on those. But C10 and C11 --- the two claims the paper's \emph{Discussion} spends the most words on (the epidemiological cgMLST tree and the anomalous non-transferability of the mcr-1 plasmid) --- are the ones we \emph{didn't} test. So the ``85\% coverage / 100\% agreement'' number is systematically biased toward the easily-testable claims. The claims we skipped are the most interesting and most contested ones in the paper.

  \item \textbf{Reference set for the backbone-conservation BLAST was hand-curated, and that biases the identity distribution upward.} We stripped the 20 esearch hits to 17 real IncX4/mcr-1/30--40 kb plasmids (dropped 3 near-duplicates); the paper's set of 26 plasmids likely includes broader IncX4 diversity (some without mcr-1, some outside the 30--40 kb window). Our 99.70--99.95\% identity is therefore \emph{expected} to be tighter than the paper's 93--98\% because our filter is stricter. We should not read our higher-identity number as ``even stronger evidence'' --- it's the same evidence on a narrower reference set. The rerun \emph{does} corroborate the qualitative claim (``highly conserved backbone''), but not more strongly than the paper.

  \item \textbf{The 100-bp chromosome length disagreement (paper 4{,}966{,}963 vs NCBI 4{,}967{,}063) was written off as ``likely paper typo''.} That is a plausible guess but we did not check the paper's supplementary materials, did not compare the paper's assembly version tag to the current NCBI \texttt{.1} record (a re-polishing or re-submission could have added 100 bp), and did not read the NCBI version history. ``Typo'' is the working assumption; the actual root cause was not investigated.

  \item \textbf{We did not check the deposited FASTA MD5 sums against any secondary NCBI mirror or the paper's supplementary FASTA (if any was submitted).} If the NCBI records have been silently updated post-publication (which does happen for polished Nanopore assemblies) our ``exact size match'' on 5 of 6 replicons could be a match to a \emph{revised} record, not the record as it stood when the paper was written.

  \item \textbf{The LLM judge was single-model + single-shot after fallback.} Argo Opus 4.7 flapped with 502s, we swapped to Argo GPT-5.2 and accepted the first response. No cross-model agreement check, no re-run with a different prompt framing. For a REPLICATED verdict on a comparatively straightforward genomic replication this is defensible, but the verdict has 1$\times$ redundancy, not the 2--3$\times$ we would want for a marginal call.

  \item \textbf{No wet-lab or expression evidence.} Every claim we tested is sequence-based. The paper's most \emph{clinically} loaded claims --- colistin MIC 8 $\mu$g/mL (Q4), non-transferability under conjugation (Q1) --- cannot be touched by this pipeline.
\end{enumerate}

Bottom line of the critique: the \emph{REPLICATED} verdict is defensible for the sequence-based
claims we tested (C1--C9), but the report should be honest that (a) we tested against curated tool
panels not against biology, (b) the two most epidemiologically-loaded claims (C10, C11) were not
touched, and (c) we validated the deposited assembly not the paper's assembly workflow.

\section{Verdict}
\label{sec:verdict}

\textbf{REPLICATED.} Every sequence-testable claim (C1--C9) independently reproduces on a live
rerun of the paper's genomic pipeline against independently-pulled NCBI records. C10 (cgMLST tree)
and C11 (non-transferability) were declared out of scope. LLM judge
\texttt{argo:gpt-5.2}: \emph{\{verdict: REPLICATED, coverage\_pct: 85, agreement\_pct: 100,
one\_line: ``NCBI assemblies confirm Ec45-2020 ST162 O55:H10; mcr-1.1 only on IncX4 33.5kb
plasmid; other plasmids/AMR largely match.'', confidence: high\}}.

The critique in \S\ref{sec:critique} does not change the verdict but does bound its
interpretation: this is ``the deposited assembly re-agrees with itself under an independent
in-silico pipeline''. It is \emph{not} ``the paper's raw reads independently re-assemble and
re-confirm the assembly''; that stronger claim would require pulling the SRA reads and re-running
Unicycler + Nanopolish, which is out of scope but achievable in a second wave.

\section{Open Questions}
\label{sec:openq}

Five open questions grounded in re-reading the paper (see \texttt{work/paper\_text.md}, extracted
from PMC10604428) plus what this replication surfaced. Full JSON in
\texttt{report/open\_questions.json}.

\paragraph{Q1. Molecular basis of non-transferability.} Why is pEc45-2020-33kb non-transferable
under standard conjugation when the same IncX4/mcr-1 backbone is transmissible in the near-identical
Chinese/European plasmids ($>$99.7\% weighted identity) it shares $\geq$96.5\% coverage with? Is it
(a) oriT loss/mutation this replication has not localized, (b) trans-inhibition by the co-resident
254 kb IncH + 190 kb IncF plasmids (Szmolka's ``conjugation competition'' speculation), or
(c) strain-level T4SS regulatory silencing? Basis: paper's own Discussion explicitly flags this
as unresolved (``the oriT was not detectable ... warrants further investigation''); our BLAST
rerun confirms the T4SS locus is essentially intact so the block is not gross sequence loss.

\paragraph{Q2. Robustness of the cgMLST cluster claim.} Is the ``Hungarian duck ST162 clusters
with Chinese human ST162 not Chinese duck ST162'' finding robust to (a) the 238-genome subsample,
(b) the K-12 MG1655 reference choice, and (c) the proprietary Ridom SeqSphere+ v9.08 defaults?
Basis: C10 is out of scope for our rerun; the paper shows a single tree with no bootstrap,
no independent scheme, and no time-calibration --- yet its ``human $\to$ duck'' direction claim
depends on this tree.

\paragraph{Q3. Mechanism + phenotype of the duplicated blaTEM-135--sul2--tet cassette.} What is
the mechanism (transposition, homologous recombination, IS-mediated capture) and evolutionary
timing (pre-existing vs recent independent acquisition) of the duplication of this cassette across
CP134087 and CP134088, and does it raise effective AMR gene copy number / MIC? Basis: paper reports
the duplication as bare fact; no flanking-IS analysis, no copy-number qPCR, no MIC-vs-single-copy
comparison. AMRFinderPlus alone cannot distinguish recent duplication from independent acquisition.

\paragraph{Q4. Modest colistin MIC (8 $\mu$g/mL) vs typical mcr-1.} Why is the colistin MIC
exactly at the EUCAST resistance breakpoint (8 $\mu$g/mL) rather than the more typical mcr-1 range
of 4--16 with mode $\sim$16? (a) Low mcr-1 expression from the IncX4 promoter in this duck strain
background, (b) compensating chromosomal mgrB/phoP/pmrAB polymorphism, or (c) fitness cost of the
two large co-resident plasmids? Basis: paper reports MIC 8 without context or chromosomal
mgrB/pmrAB screening; our AMRFinderPlus rerun does not comprehensively cover the E. coli colistin
core-mutation panel. The modest MIC materially affects the One-Health interpretation.

\paragraph{Q5. Hidden gene content of the ``AMR-free'' 101 kb + 5 kb plasmids.} Both the paper
and this replication dismissed CP134086 and CP134090 as ``no AMR/VG determinants'' --- but
neither actually annotated their ORFs. Col156 plasmids commonly carry colicins; p0111 plasmids
are implicated in APEC virulence in the literature. Basis: our AMRFinderPlus rerun only calls
curated AMR/VF panels. ``No hits from AMRFinderPlus'' $\neq$ ``no biologically meaningful genes''.
This is exactly the kind of ``data not shown'' gap Rick flagged.

Concrete next steps for each Q are enumerated in \texttt{report/open\_questions.json}.

\section{Companion Artifacts}
\label{sec:companion}

\begin{itemize}
  \item \texttt{report/REPORT.md} --- original markdown report (kept per standard).
  \item \texttt{report/open\_questions.json} --- machine-readable version of \S\ref{sec:openq}.
  \item \texttt{report/workflow.md} --- workflow narrative + tool inventory + effort estimate.
  \item \texttt{report/artifacts\_summary.md} --- inventory + traces.
  \item \texttt{report/failure\_analysis.md} --- honest failure/gap analysis.
  \item \texttt{report/evidence/} --- amrfinder\_full.tsv, mlst\_warwick.tsv, mlst\_pasteur.tsv, plasmidfinder.tsv, serotype.tsv, blast\_incx4\_v1.tsv, blast\_incx4\_v2.tsv, blast\_incx4\_v2\_summary.json, llm\_judge\_verdict.json, llm\_judge\_model.txt, amrfinder.log.
  \item \texttt{work/} --- paper\_pmc.xml, paper\_text.md, pmc\_link.json, pubmed\_meta.json, mcr1\_ref.faa, blast\_summary.py, llm\_judge\_prompt.txt, genomes/.
  \item \texttt{paper.pdf} --- backfilled from MDPI open-access (CC BY 4.0) during 2026-07-05 8-artifact pass.
  \item \texttt{extraction/marker.md} --- pdftotext fallback extraction; central Marker corpus copy pending sha256 resolution.
  \item \texttt{extraction/nougat.mmd} --- pending central Nougat parse (GPU-only; stub with sha256 header written for later corpus sweep).
\end{itemize}

\end{document}
